\section{格在语义、约束与 Galois 连接中的角色}
\label{sec:04-Discrete-Mathematics/08-Rings-and-Lattices/03}

一个静态分析工具要在不运行程序的前提下判断 \texttt{a / b} 会不会除零。\texttt{b} 的取值来自此前所有分支、循环与输入，把它们枚举一遍是不可能的。工具的做法是干脆不记具体数值，只记一个标签：\texttt{b} 是正、是负、是零，还是说不准。于是每个变量的状态从一个无穷集合塌缩成四个符号之一，计算量降到可以承受。

这一步交换了什么，要说清楚。把具体的取值集合换成标签，是在放松：\(\{2,4,6\}\) 与 \(\{7\}\) 都记成``正''，信息丢了一部分。反过来，拿到``正''这个标签要回答关于程序的问题时，只能把它读成全体正数——不敢再窄，否则结论就不可靠了。一边放松、一边收紧，两个方向的映射配成一对。上一节留下的单调映射与最小不动点，正是在这样一对映射的框架里变成可执行的算法。

\subsection{一对反向映射，配对条件是什么}

设 \((P,\le)\) 与 \((Q,\sqsubseteq)\) 是两个偏序集，\(\alpha:P\to Q\) 负责抽象化，\(\gamma:Q\to P\) 负责具体化。称 \((\alpha,\gamma)\) 是 Galois 连接，若对一切 \(p,q\)

\[
\alpha(p)\sqsubseteq q
\quad\Longleftrightarrow\quad
p\le\gamma(q).
\]

这一行读起来抽象，意思却很实在：\(q\) 是 \(p\) 的一个合法描述，当且仅当 \(q\) 允许的全部具体值都包含 \(p\)。左边站在抽象域里比较，右边站在具体域里比较，条件说这两种检查等价，所以你可以任选一边做，另一边自动成立。附带的结论是 \(\alpha(p)\) 必为 \(p\) 的最佳抽象——所有合法描述里最紧的那个。

\begin{center}
\begin{tikzpicture}[x=1cm,y=1cm,font=\footnotesize,>=stealth]
  \draw[black!45,fill=blue!5] (0,0) rectangle (5.4,5.0);
  \node[black!75,font=\small] at (2.7,5.4) {具体域 \(P\)（按包含排序）};
  \draw[black!45,fill=orange!7] (9.0,0) rectangle (12.6,5.0);
  \node[black!75,font=\small] at (10.8,5.4) {抽象域 \(Q\)};

  \fill (1.4,1.0) circle (2pt);
  \node[right,black!80] at (1.6,1.0) {\(p=\set{2,4,6}\)};
  \fill (1.4,3.4) circle (2pt);
  \node[right,black!80] at (1.6,3.4) {\(\gamma(\alpha(p))=\set{1,2,3,\dots}\)};
  \draw[->,dashed,black!60] (1.4,1.25) -- (1.4,3.15);
  \node[left,black!65] at (1.3,2.2) {闭包};

  \fill (10.8,3.4) circle (2pt);
  \node[right,black!80] at (11.0,3.4) {正};
  \fill (10.3,1.2) circle (2pt);
  \node[right,black!80] at (10.5,1.2) {零};
  \fill (10.8,4.4) circle (2pt);
  \node[right,black!80] at (11.0,4.4) {\(\top\)};

  \draw[->,thick,blue!60] (2.9,1.15) .. controls (6.0,1.9) and (7.6,2.7) .. (10.55,3.3);
  \node[black!75,above] at (6.6,2.55) {\(\alpha\)：放松，记住方向就够};
  \draw[->,thick,orange!70] (10.6,3.55) .. controls (7.6,4.6) and (5.0,4.4) .. (1.6,3.6);
  \node[black!75,below] at (6.3,4.15) {\(\gamma\)：收紧，读成它允许的全部值};
\end{tikzpicture}

\emph{一次往返把 \(p\) 抬到了 \(\gamma(\alpha(p))\)，位置只升不降，抬上去的部分就是这次抽象付出的精度。再走一遍往返不会继续上升——闭包算子是幂等的，损失一次性发生完。}
\end{center}

复合 \(\gamma\circ\alpha\) 单调、幂等，而且满足 \(p\le\gamma(\alpha(p))\)，这三条合起来就是闭包算子的定义；另一头的 \(\alpha\circ\gamma\) 则是收缩的核算子。幂等性在实践中直接可用：同一层抽象反复做没有意义，想要更准只能换一个更细的抽象域，不能靠多做几遍。

\begin{example}
把 Galois 连接放在``对象''与``性质''之间：一组对象映到它们共有的全部性质，一组性质映到满足这些性质的全部对象。这对映射自动构成 Galois 连接，而闭包算子的不动点对——一组对象连同恰好刻画它们的那组性质——就是形式概念分析里的概念。三角形按``等边、等腰、直角''归类时，从``等边''出发走一个往返会回到``等边且等腰''，因为等边三角形全都等腰。闭包算子把一组性质补全成它蕴含的全部性质，概念格的层次就是这样被算出来的。
\end{example}

\subsection{可靠性由哪一半负责}

抽象解释的承诺是单向的：分析器可以多报，不能漏报。可能除零的地方它必须报出来，不可能除零的地方它报了也只是浪费人力。这条承诺的技术名称是可靠性（soundness），它的来源是一条覆盖条件——抽象域上的转移函数必须罩住具体域上的真实转移，即先做一步具体计算再抽象，不超过先抽象再做一步抽象计算。有了它，抽象格上算出的结果经 \(\gamma\) 还原后一定包住真实行为。

精度则是另一笔账，由抽象域的选择决定。符号域太粗，``正''减``正''只能得到 \(\top\)，循环里几步之后所有变量都变成``说不准''。换成区间域，记住 \(x\in[0,100]\) 这种信息，能过的检查明显变多，代价是每步转移的计算贵得多；再往上还有关系型的域，能表达 \(x\le y\) 这类变量之间的约束。粗与细之间没有免费的一档。

同一套机制在约束求解里叫传播。求解器给每个变量维护一个候选值域，每条约束都是一个单调映射：把与当前候选相矛盾的取值删掉。值域按包含关系构成格，删值就是沿格向下走，走到不再有值可删为止。弧一致性算法算的正是这个不动点。它同样不承诺解出问题——传播停下来时可能还剩多个候选，那时才轮到搜索上场。

\subsection{分配律塌了，逻辑跟着变}

上一节说过子空间格不分配，这里把后果讲完。\(\R^2\) 里取三条互不相同的过原点直线 \(V,W_1,W_2\)。它们两两相交只有原点，所以 \((V\cap W_1)+(V\cap W_2)\) 是零空间；而 \(W_1+W_2\) 已经是整个平面，\(V\cap(W_1+W_2)=V\) 非零。分配律的两端差了一整条直线。

这不是一个技术性的小毛病。量子逻辑用子空间格代替布尔代数来描述命题，直接后果就是``粒子穿过左缝或右缝''可以为真，而``穿过左缝''与``穿过右缝''都不为真——在布尔代数里这是矛盾，在子空间格里只是分配律不成立。经典逻辑的幂集结构、量子命题的子空间结构、约束系统的值域结构是三种不同的格，把其中一种的推理规则搬到另一种上，得到的是伪证明。所以读到一步``展开''或者一步``排中''时，值得停一下：这一步用的是格的哪条性质，当前这个格有没有。

\subsection{迭代为什么一定停得下来}

上一节的 Knaster--Tarski 定理保证了最小不动点存在，但存在不等于算得到。真正让分析器能交差的是格的高度——从底到顶最长的严格上升链有多长。高度有限时，从 \(\bot\) 出发反复施加单调转移并取 join，每一步要么严格上升要么已经稳定，链长封顶，迭代必在有限步内停住。

\begin{center}
\begin{tikzpicture}[x=1cm,y=1cm,font=\footnotesize,>=stealth]
  \draw[black!55] (0,0) -- (-1.3,1.3);
  \draw[black!55] (0,0) -- (0,1.3);
  \draw[black!55] (0,0) -- (1.3,1.3);
  \draw[black!55] (-1.3,1.3) -- (0,2.6);
  \draw[black!55] (0,1.3) -- (0,2.6);
  \draw[black!55] (1.3,1.3) -- (0,2.6);
  \fill[white] (0,0) circle (0.3);
  \draw[black!60] (0,0) circle (0.3);
  \node at (0,0) {\(\bot\)};
  \fill[white] (-1.3,1.3) circle (0.3);
  \draw[black!60] (-1.3,1.3) circle (0.3);
  \node at (-1.3,1.3) {负};
  \fill[white] (0,1.3) circle (0.3);
  \draw[black!60] (0,1.3) circle (0.3);
  \node at (0,1.3) {零};
  \fill[white] (1.3,1.3) circle (0.3);
  \draw[black!60] (1.3,1.3) circle (0.3);
  \node at (1.3,1.3) {正};
  \fill[white] (0,2.6) circle (0.3);
  \draw[black!60] (0,2.6) circle (0.3);
  \node at (0,2.6) {\(\top\)};
  \node[black!70] at (0,-0.8) {符号格：高度 3};

  \draw[black!30] (3.6,-0.4) -- (3.6,3.6);

  \foreach \k in {0,1,2,3}{
    \fill (5.6,0.55*\k) circle (1.8pt);
  }
  \draw[->,black!60] (5.6,0.1) -- (5.6,0.45);
  \draw[->,black!60] (5.6,0.65) -- (5.6,1.0);
  \draw[->,black!60] (5.6,1.2) -- (5.6,1.55);
  \node[right,black!80] at (5.8,0.0) {\([0,0]\)};
  \node[right,black!80] at (5.8,0.55) {\([0,1]\)};
  \node[right,black!80] at (5.8,1.1) {\([0,2]\)};
  \node[right,black!80] at (5.8,1.65) {\([0,3]\)};
  \draw[black!60,dashed] (5.6,1.85) -- (5.6,2.4);
  \node[right,black!60] at (5.8,2.2) {没有尽头};
  \fill[red!60] (5.6,3.2) circle (2pt);
  \node[right,black!80] at (5.8,3.2) {\([0,+\infty)\)};
  \draw[->,thick,red!60] (5.3,1.7) .. controls (4.2,2.4) and (4.4,3.0) .. (5.3,3.2);
  \node[left,red!70] at (4.2,2.5) {widening};
  \node[black!70] at (6.4,-0.8) {区间格：高度无限};
\end{tikzpicture}

\emph{左边的符号格四层封顶，任何单调迭代最多爬三步。右边的区间格可以无限上升，循环变量每转一圈把上界推大 1，迭代永不停止；widening 强行把这条链跳到顶，用一次大幅度的精度损失换来终止。}
\end{center}

widening 的取舍值得单独说一句。它不是工程上的将就，而是这类分析的固有代价：既要在有限时间内给出答案，又要覆盖无限多的执行路径，总得在某处主动放弃精度。数值分析里步长与误差的权衡、统计里偏差与方差的权衡，形状和它一样——三处都是同一个交易在不同场合出现。

\subsection{这一章留下什么}

格论在本书里不追求自身完备，它提供的是三个可以随身带走的检查。第一，你写的那个合并函数当真是 join 吗？交换、结合、幂等缺一条，副本收敛与迭代终止的保证就不成立。第二，你手上那对往返映射构成 Galois 连接吗？构成，则往返一次的损失有界且幂等，抽象结果可以放心具体化；不构成，则``近似''二字没有兑现的依据。第三，你的推理暗中用了分配律或补元吗？子空间格的例子说明这两条随时会不在。

回头看整章，两条路线的形状其实一致。环那边是往集合上叠运算，叠一条要求换一种可以放心做的操作，去掉``非零元可逆''就掉出零因子，除法随之失效；格这边是从既有的次序里读出运算，多要求一条完备性就换来不动点的存在，多要求一条有限高度就换来算法的终止。写下``这个结构满足什么''与追问``所以我能做什么''，始终是同一个问题的两面。下一章回到群那条线，把群搬到矩阵上去看，用线性代数的工具把对称拆成不可再分的零件。
